\FloatBarrier
\section{Modelling of cement hydration and its reliance on SEM images}

% aus dem ct/fib-sem-antrag
As the demands on building materials increase, so do the requirements for characterising and modelling the underlying material systems.
There are advanced approaches to modelling and simulating the strength, shape retention and durability characteristics of concretes \cite{bishnoi.2009, kumar.2017, zhou.2019, ouzia.2019}.
Due to the high complexity, multi-scale heterogeneity and temporal changes, this is often only possible in a highly simplified form and in a very limited scale.

The microstructure of concretes is simulated in a variety of models at different scales \cite{bishnoi.2009, kumar.2017, zhou.2019, ouzia.2019}.
The mesostructure (between atomistic and micro scale) of e.g. \ac{CSH} phases and the porosity was also modelled \cite{ioannidou.2016, yu.2016}.
Atomistic modelling is used to fundamentally simulate the structural composition \cite{kumar.2017,mishra.2017,kunhimohamed.2020,casar.2023}.
Finally, there are modern micromechanical models that attempt to close the transition from meso- to macromodelling \cite{bernard.2003, hlobil.2016, gobel.2018}.

The hydration of cement over time in 3D using the sheet growth model using a level set based method \cite{nguyen-tuan.2019,nguyen-tuan.2023} (see Figure~\ref{fig:sheet_growth}).

% from a talk of long from the ICCC 2019, slide 14/18 but not in the paper
\begin{figure*}[htb]
    \centering

    \ifdefined\PRINTABLE
        \begin{subfigure}[t]{0.49\linewidth}
                \includegraphics[width=\textwidth]{literature/sheet_grow_anim/sg01.jpg}
        \end{subfigure}%
        \hfill%
        \begin{subfigure}[t]{0.49\linewidth}
            \includegraphics[width=\textwidth]{literature/sheet_grow_anim/sg04.jpg}
        \end{subfigure}
    \else
        \animategraphics[palindrome,autoplay,width=.49\linewidth]{3}{literature/sheet_grow_anim/sg}{01}{04}
    \fi

    \caption{
        Visualisation of the growth of \ac{CSH} in a representative domain of \num{.5}\,$\times$\,\num{.5}\,$\times$\,\SI{.5}{\micro\meter\cubed} between \SI{15}{\minute} and \SI{4}{\hour} , based on the sheet gowth model \cite{nguyen-tuan.2019,nguyen-tuan.2023}.% in the second paper 300x300x2200nm is used due to some restrictions.
    }
    \label{fig:sheet_growth}
\end{figure*}

For larger domains, \textcite{nguyen-tuan.2024} utilized the sheet growth model allowing to simulate the growth of low and high density \ac{CSH} in 2D based on a segmented image of polished alite after \SI{1}{\day} of hydration (see Figure~\ref{fig:level_set}).

\begin{figure*}[htb]
    \centering

    \begin{subfigure}[t]{0.32\linewidth}
        \includegraphics[width=\textwidth]{literature/level_set_anim/ls01.jpg}
    \end{subfigure}%
    \hfill%
    \begin{subfigure}[t]{0.32\linewidth}
        \ifdefined\PRINTABLE
            \includegraphics[width=\textwidth]{literature/level_set_anim/ls04.jpg}
        \else
            \animategraphics[palindrome,autoplay,width=\textwidth]{6}{literature/level_set_anim/ls}{01}{09}
        \fi
    \end{subfigure}%
    \hfill%
    \begin{subfigure}[t]{0.32\linewidth}
        \includegraphics[width=\textwidth]{literature/level_set_anim/ls07.jpg}
    \end{subfigure}

    \caption{Visualisation of the growth of \ac{CSH} (blue) in the low and high density variant on top of \ac{C3S} (yellow) by \textcite{nguyen-tuan.2024}, based on the level set method.}
    \label{fig:level_set}
\end{figure*}

However, the basis for all realistic modelling is a quantitative experimental characterisation of the corresponding structures \cite{gobel.2018a}.

{\color{red} % erstmal nur so dahin gequatscht
\textcite{hlobil.2022} crated a model called \textsc{CEMHYD3D} which was compared to volumetric \ac{mCT} data of differently aged cement.
Models like those (\cite{hlobil.2022, nguyen-tuan.2024}) allow to simulate the development of the phase composition over time which can be validated using various measurements.
3D-based simulations inherently have a better chance to closely image the reality, while 2D-simulations, as image sections are always missing a lot of information of the grain.
However, they are much faster and allow to simulate larger areas reflecting a more representative are, which is very important due to the inhomogeneity of cement (see Section~\ref{sec:2D-representativeness})
% simulationspaper der hydratation \cite{hlobil.2022}
% CEMHYD3D -> Fig 3 anteile der phasen vgl mit Messungen
% Literatur aus Long übernehmen und Modelle erwähnen
}

% Handlungsanweise - wie kann man REM mit Methoden kombinieren um Baustoffe zu beschreiben - welche Techniken gibt es?
% Einfache Tools, weil sonst machts keener